\subsection{Exploratory extension: more baselines, longer training, per-$k$ tuning}
\label{sec:extension}
These analyses were added after the main PopSign results and the audit that followed (\texttt{PREREG\_ADDENDUM.md}, Part B); they are exploratory.
\begin{table}[t]
\centering\small
\begin{tabular}{lrrrrrr}
\toprule
Method & 1 & 2 & 5 & 10 & 20 & all \\
\midrule
DTW 1-NN (main) & 17.2 & 22.1 & 31.0 & 38.6 & 45.9 & 63.4 \\
Euclidean 1-NN (main) & 18.9 & 24.7 & 35.0 & 42.7 & 47.5 & 63.2 \\
DTW k-NN (kappa on validation) & 17.2 & 22.1 & 31.0 & 38.6 & 45.9 & 63.4 \\
Fisher-weighted DTW & 17.2 & 22.4 & 30.3 & 36.6 & 43.5 & 59.2 \\
MiniRocket + ridge & 11.6 & 12.5 & 30.0 & 40.0 & 50.7 & 70.8 \\
DTW, tuned per k & 17.2 & 22.1 & 31.0 & 38.6 & 45.9 & 63.4 \\
Soft-DTW prototype & 19.7 & 18.7 & 15.7 & 14.4 & -- & -- \\
Embedding (trained on other signs) + 1-NN & 47.8 & 54.5 & 62.0 & 64.9 & 68.1 & 75.3 \\
1D-CNN (main) & 23.2 & 32.0 & 48.2 & 61.1 & 71.6 & 85.5 \\
1D-CNN, tuned per k & 23.2 & 32.0 & 48.3 & 61.1 & 71.6 & 85.1 \\
1D-CNN, 200 epochs & 23.2 & 32.1 & 49.1 & 61.9 & 72.8 & 87.9 \\
BiGRU (main) & 20.0 & 28.7 & 43.9 & 56.7 & 67.9 & 84.9 \\
BiGRU, 200 epochs & 20.0 & 30.1 & 45.1 & 58.3 & 67.5 & 86.6 \\
Transformer (main) & 23.7 & 30.3 & 45.9 & 57.0 & 66.2 & 85.9 \\
Transformer, 200 epochs & 23.7 & 30.1 & 45.2 & 57.2 & 66.8 & 86.5 \\
\bottomrule
\end{tabular}

\caption{PopSign, 50 signs: top-1 accuracy (\%) of additional methods and variants (exploratory) beside the main ones.}
\label{tab:ext}
\end{table}
\paragraph{Other nearest-neighbour variants.} Voting among several nearest references did not help: the number of neighbours chosen on validation was
one for both DTW and Euclidean. Weighting feature dimensions by a Fisher score estimated from the references did not help either
(\accWdtwVFiftyKAll\% with all examples against \accDtwVFiftyKAll\% for DTW).
\paragraph{A stronger time-series baseline.} MiniRocket with a ridge classifier \citep{dempster2020rocket} is poor with one or two examples per sign
(\accRocketVFiftyKOne\% and \accRocketVFiftyKTwo\%, against \accDtwVFiftyKOne\% and \accDtwVFiftyKTwo\% for DTW), about level with DTW at $k=5$
and 10 (\accRocketVFiftyKFive\% and \accRocketVFiftyKTen\% against \accDtwVFiftyKFive\% and \accDtwVFiftyKTen\%), and above it from
$k=20$ (\accRocketVFiftyKTwenty\% against \accDtwVFiftyKTwenty\%, and \accRocketVFiftyKAll\% against \accDtwVFiftyKAll\% with all examples)
but well below the neural models (\accCnnVFiftyKAll\% for the CNN).
\paragraph{Tuning per $k$ and longer training.} Re-selecting every method's settings separately for $k\in\{1,5,10,\text{all}\}$ on validation
changed no cell by more than \perkMaxChange{} points. Training the learned models for 200 epochs with a step cap of 12{,}000 instead of the
main budget changed accuracy by between \longMinGainCnn{} and \longMaxGainCnn{} points for the CNN, \longMinGainGru{} and \longMaxGainGru{} for
the BiGRU and \longMinGainTrans{} and \longMaxGainTrans{} for the Transformer across all $k$ (Table~\ref{tab:ext}), so the main results
were not far from converged and the gap to DTW did not close.

\subsection{Descriptive analyses of why (exploratory)}
We looked at three possible explanations; none clearly accounts for the gap. (1)~\emph{Hand slot.} MediaPipe assigns each hand to a left or right slot, and test signers differ:
each signer uses predominantly one slot (at least \whyMinDominantShare\% of hand-present frames), with the share showing only the left slot between \whyLeftOnlyMin\% and
\whyLeftOnlyMax\% and only the right slot between \whyRightOnlyMin\% and \whyRightOnlyMax\% across signers; \whyNLeftMajority{} of 8 mostly use the left
slot. Swapping the slots of the test queries lowers accuracy only modestly (DTW with all examples \whySwapDtwNormalAll\% to \whySwapDtwSwappedAll\%,
CNN \whySwapCnnNormalAll\% to \whySwapCnnSwappedAll\%), so the models are not tied to a single slot. The large gain from putting the more active hand in the first slot
(Table~\ref{tab:abl}) is consistent with pooling a sign's few examples across signers who use different hands, which we did not test directly.
(2)~\emph{Movement.} Per-sign wrist path length has no clear relation to the per-sign accuracy gap between the CNN and DTW
(Spearman $\rho=\whyMotionRho$, 95\% bootstrap interval \whyMotionLo{} to \whyMotionHi{} over 50 signs).
(3)~\emph{Feature scale or preprocessing.} Both were tuned for DTW on validation (Section~\ref{sec:results}). These analyses are descriptive; a causal account remains open.
